\section{Introducción}
\label{sec:introduccion}

% ==== REEMPLAZA este texto por la introducción de TU caso ====
Muchos problemas de ingeniería —circuitos eléctricos, estructuras, redes de tuberías,
balances de masa y procesos térmicos— se modelan como un \textbf{sistema de ecuaciones
algebraicas lineales} de la forma

\begin{equation}
\mathbf{A}\,\mathbf{x} = \mathbf{b}
\label{eq:sistema}
\end{equation}

donde $\mathbf{A}\in\mathbb{R}^{n\times n}$ es la matriz de coeficientes,
$\mathbf{x}\in\mathbb{R}^{n}$ es el vector de incógnitas y
$\mathbf{b}\in\mathbb{R}^{n}$ es el vector de términos independientes.

El objetivo de este reporte es resolver el caso de estudio asignado mediante métodos
numéricos \textbf{directos} (Eliminación Gaussiana con pivoteo, Descomposición LU) e
\textbf{iterativos} (Gauss--Seidel), así como algoritmos para \textbf{matrices especiales}
(Thomas y Cholesky). Se compara la implementación ``a mano'' contra los \emph{solvers}
profesionales de \texttt{NumPy}/\texttt{SciPy} y se somete el sistema a una simulación de
\textbf{Monte Carlo} para evaluar su \textbf{condicionamiento} y sensibilidad ante
incertidumbre en los datos.

% Explica aquí: qué fenómeno físico modela tu caso, por qué es un sistema lineal,
% y qué decisión de ingeniería depende de la solución.
